% !TeX root = ../../Book3.tex
% 纲要标题：## 附录 F　局部标准基、幂级数与切锥计算

\chapter{局部标准基、幂级数与切锥计算}
\label{b3:app:F}
\BookChapterNumber{b3:app:F}

% 纲要细目（与计划纲要同步）：
% > 所需基础是第2、8、11、13章的局部化、全局 Gröbner 基、滤过与完备化。正文中的维数、深度和局部环结构不以本附录为前提；局部计算专题读者可在完备化后插读。


\begin{chapterabstract}
局部首项约化可能不断产生更高次项。Mora 弱正规形的单位乘子与暂存余项机制解释了如何得到有限证书，本附录通过算例核验这一机制，并证明局部次数序下的形式除法。初始齐次理想与伴随分次环用于切锥和 Hilbert--Samuel 数据，标准基及截断证书还用于 Jacobian 商等有限长度计算。
\end{chapterabstract}

\begin{learninggoals}
\begin{itemize}
\item 区分全局序、局部序、混合序以及它们所确定的工作环。
\item 用有限等式和单位条件验证弱正规形证书，并区分 Mora 方法的基本机制、有限步终止与形式收敛。
\item 从全部初始关系计算切锥，辨认任意生成元的最低次项可能遗漏的关系。
\item 给出有限长度和重数计算所需的截断依据，用生成元变换和核关系验证局部模计算。
\end{itemize}
\end{learninggoals}

在原点的局部环中，\(1-x\) 是单位，所以 \(x\) 与 \(x-x^2=x(1-x)\) 生成同一个理想。若按局部序把低次项 \(x\) 作为首项，却仍照搬全局除法中逐次消去首项的做法，用 \(x-x^2\) 反复约化 \(x\)，便先余 \(x^2\)，再余 \(x^3\)，如此继续而不会停。这不是理想成员资格失效，而是所用的约化过程没有利用局部单位。

\ProseParagraphBreak

局部标准基必须允许在等式中出现适当的单位乘子，使“属于理想”的证明仍是有限的，同时控制余式所采用的局部单项式次序。形式除法还要在完备环里解释无限尾项。把这些计算与伴随分次环相连，便可从原点附近的最低次项读取切锥、局部长度及某些奇点不变量。

\ProseParagraphBreak

本附录需要\chapref{b3:ch:02}的局部化、\chapref{b3:ch:08}的全局 Gröbner 基、\chapref{b3:ch:11}的滤过与 Hilbert--Samuel 理论，以及\chapref{b3:ch:13}的完备化。最后讨论自由消解时调用\chapref{b3:ch:21}。前三节的具体计算以完备化为止。

\ProseParagraphBreak
记 \(x=(x_1,\ldots,x_n)\)，\(k\) 为域。\cref{b3:tab:F-computational-settings}比较三种工作环中的单位和计算完成方式。全局与局部两行给出有限多项式等式，形式行中的商系数一般仍是无限对象，须按各阶系数检查相容性，或另给这些系数的有限表示。

\begin{table}[htbp]
\centering
\caption{全局、局部与形式计算的完成方式}
\label{b3:tab:F-computational-settings}
\begin{tabularx}{\textwidth}{@{}>{\raggedright\arraybackslash}p{.17\textwidth}>{\raggedright\arraybackslash}p{.20\textwidth}>{\raggedright\arraybackslash}p{.25\textwidth}>{\raggedright\arraybackslash}X@{}}
\toprule
工作环 & 方法与单位 & 完成方式 & 成员关系的依据 \\
\midrule
\(k[x]\) & 全局 Gröbner；单位为 \(k^\times\) & 良序支持有限步停止 & 有限等式 \(f=\sum_i a_i g_i\)，\(a_i\in k[x]\) \\
\(k[x]_{(x)}\) & 局部 Mora；单位分子的常数项非零 & 带单位乘子的有限等式 & \(uf=\sum_i a_i g_i\)，\(u(0)\ne0\)，\(u,a_i\in k[x]\) \\
\(k[\![x]\!]\) & 局部次数序除法；单位级数常数项非零 & 各阶系数稳定后取完备极限 & 形式恒等式 \(f=\sum_i q_i g_i\)；\(q_i\) 通常仍是无限级数 \\
\bottomrule
\end{tabularx}
\end{table}

% !TeX root = ../../Book3.tex
% 纲要标题：### F.1 局部化、局部序与标准基

\section{局部化、局部序与标准基}
\label{b3:sec:F01}

在一个点附近，低阶项比高阶项更能反映局部结构。改变单项式次序后，哪些多项式能够当作单位也会改变，因此必须同时确定计算所在的局部化环。

% 纲要细目（与计划纲要同步）：
% * 比较全局、局部及混合序；正式算例准确识别混合局部化的单位与非单位，不把混合序都当作点局部环
% * 以首单项式理想相等定义标准基，证明分式首项相容与有限首单项式生成；存在性与有效求基分开
% * 完整证明有限阶梯给出局部长度；用同一理想比较局部最小生成元数、首单项式生成元数和阶梯长度
% #### F.1.1 全局序、局部序、混合序与相应局部化
% #### F.1.2 首项理想与标准基的定义
% #### F.1.3 标准基、标准单项式与局部最小生成元

\subsection{全局序、局部序、混合序与相应局部化}
\label{b3:subsec:F01:01}

在多项式环中，常把高次项放在前面；在一个点附近，最低次非零项往往更重要。它决定元素在极大理想的哪一层出现，并给出切锥中的齐次关系。局部项序正是为了先处理这些低阶项。不过，换序的同时必须说明工作环和除法方法也发生了什么变化。

\begin{definition}\label{b3:def:F-monomial-orders}
设 \(k\) 为域，\(n\ge1\)，\(S=k[x_1,\ldots,x_n]\)。若单项式上的全序 \(>\) 与乘法相容，则称它为单项式次序。若每个 \(x_i>1\)，则称 \(>\) 为全局序；若每个 \(x_i<1\)，则称 \(>\) 为局部序；若部分变量大于 \(1\)、部分小于 \(1\)，则称 \(>\) 为混合序。对非零多项式 \(f\in S\)，取其有限个非零项中最大的单项式为 \(\operatorname{LM}_{>}(f)\)，相应的带系数项为 \(\operatorname{LT}_{>}(f)\)。
\end{definition}

第8章的“单项式序”默认要求良序；本附录使用“单项式次序”容纳局部与混合情形。若每个变量都大于 \(1\)，乘法相容性与 Dickson 引理仍保证良序，便回到第8章的定义。例如先将总次数较小的单项式规定为较大，同次再用固定字典序比较，就得到局部次数序。在此序中，\(x-y^2\) 的首项是 \(x\)。这里的大小方向与全局次数序相反：局部次数序先取低次项，全局次数序先取高次项；这一比较只针对次数序，不是任意全局序或局部序的定义。

\ProseParagraphBreak
局部序有下降链
\[
1>x>x^2>x^3>\cdots.
\]
沿此链首单项式越来越小，总次数却越来越大；在局部次数序中，同次单项式之间也可以下降。因此，“首项严格下降”不再保证有限步结束，全局多项式除法所用的良序终止证明在这里失效。

\begin{definition}\label{b3:def:F-order-localization}
设 \(k\) 为域，\(S=k[x_1,\ldots,x_n]\)，\(>\) 是其单项式次序。对非零 \(u\in S\)，记 \(\operatorname{LC}_{>}(u)\) 为首项 \(\operatorname{LT}_{>}(u)\) 的系数。令
\[
U_{>}=\{u\in S:\operatorname{LM}_{>}(u)=1\},
\qquad S_{>}=U_{>}^{-1}S.
\]
称 \(S_{>}\) 为该次序相应的局部化。对 \(f/u\ne0\)，定义
\[
\operatorname{LM}_{>}(f/u)=\operatorname{LM}_{>}(f),
\qquad
\operatorname{LT}_{>}(f/u)
=\operatorname{LT}_{>}(f)/\operatorname{LC}_{>}(u).
\]
\end{definition}
\symidx{\(U_>\)、\(S_>\)：单项式次序确定的乘法集与局部化}

乘法相容性给出非零多项式的公式
\[
\operatorname{LM}_{>}(fg)
=\operatorname{LM}_{>}(f)\operatorname{LM}_{>}(g),
\qquad
\operatorname{LT}_{>}(fg)
=\operatorname{LT}_{>}(f)\operatorname{LT}_{>}(g).
\]
因此 \(U_{>}\) 是乘性集。若 \(f/u=g/v\)，则 \(fv=gu\)；利用 \(\operatorname{LM}_{>}(u)=\operatorname{LM}_{>}(v)=1\) 比较两边首项，得到
\(\operatorname{LM}_{>}(f)=\operatorname{LM}_{>}(g)\) 及
\(\operatorname{LT}_{>}(f)\operatorname{LC}_{>}(v)=\operatorname{LT}_{>}(g)\operatorname{LC}_{>}(u)\)。这说明上述首单项式与首项都不依赖分式表示，乘法公式也随之延伸到 \(S_{>}\)。

\ProseParagraphBreak
全局序时 \(U_{>}=k^\times\)，所以 \(S_{>}=S\)；局部序时 \(U_{>}\) 是常数项非零的多项式，故
\[
S_{>}=k[x_1,\ldots,x_n]_{(x_1,\ldots,x_n)}.
\]
混合序所对应的环要由具体次序计算。例如先按 \(x\) 次数作全局比较，同次再以 \(1>y>y^2>\cdots\) 比较，则 \(U_{>}=k[y]\setminus(y)\)，相应环为 \(k[y]_{(y)}[x]\)。在这个指定的块次序中，\(1-y\) 成为单位，\(1-x\) 仍不可逆，完整核验见\cref{b3:exm:F-mixed-order}；任意混合序的工作环不能一概写成这个环。

\ProseParagraphBreak
本附录按次序选择 \(U_{>}\) 及 \(S_{>}\)，所以从全局序改为局部序或混合序时，还须重新确定工作环及其单位。原理想的生成等式、成员关系和标准基条件，都应在新工作环中解释。

\subsection{首项理想与标准基的定义}
\label{b3:subsec:F01:02}

\begin{definition}\label{b3:def:F-standard-basis}
设 \(k\) 为域，\(S=k[x_1,\ldots,x_n]\)，\(>\) 是其单项式次序，\(R=S_{>}\)，\(I\subset R\) 为理想。其首项理想指多项式环 \(S\) 中的单项式理想
\[
L_{>}(I)=(\operatorname{LM}_{>}(f):0\ne f\in I)\subset S.
\]
若 \(G=(g_1,\ldots,g_s)\) 是 \(I\) 的一个由非零元素组成的有限生成组，并满足
\[
L_{>}(I)=(\operatorname{LM}_{>}(g_1),\ldots,
\operatorname{LM}_{>}(g_s)),
\]
则称 \(G\) 为 \(I\) 关于 \(>\) 的\term{标准基}[standard basis]；零理想取空列作为标准基。本附录采用包含全局 Gröbner 基的这一约定。
\end{definition}
\symidx{\(L_>(I)\)：局部或混合次序下的首项理想}

把首项理想放在 \(S\) 中，可以用同一个指数集合记录被首单项式整除的单项式；原理想及其生成关系则始终在 \(R\) 中理解。由于非零首系数在域 \(k\) 中可逆，用 \(\operatorname{LT}_{>}\) 代替 \(\operatorname{LM}_{>}\) 也生成同一个理想。下面的有限性结论只证明标准基存在，尚未给出从指定生成组计算标准基的算法。

\begin{proposition}\label{b3:prop:F-standard-existence}
设 \(k\) 为域，\(S=k[x_1,\ldots,x_n]\)，\(>\) 为乘法相容的单项式全序，\(R=S_{>}\) 为相应局部化。则每个 \(R\) 的理想都有有限标准基。
\end{proposition}
\begin{proof}
零理想取空列即可。设 \(I\ne0\)。多项式环 \(S\) 是 Noether 环，故单项式理想 \(L_{>}(I)\) 可由有限多个元素的首单项式生成；为这些单项式选择原像 \(g_1,\ldots,g_s\in I\)。环 \(R\) 也是 Noether 环，另取 \(I\) 的一个由非零元素组成的有限生成组并加入该列。所得列生成 \(I\)，其首单项式生成 \(L_{>}(I)\)。

对列中的每个元素写 \(g_i=f_i/u_i\)，其中 \(f_i\in S\)、\(u_i\in U_{>}\)。因为 \(f_i=u_i g_i\) 且 \(u_i\) 在 \(R\) 中可逆，以 \(f_i\) 代替 \(g_i\) 不改变在 \(R\) 中生成的理想。首单项式的乘法公式又给出
\[
\operatorname{LM}_{>}(f_i)
=\operatorname{LM}_{>}(u_i)\operatorname{LM}_{>}(g_i)
=\operatorname{LM}_{>}(g_i).
\]
因此逐个清除分母后得到由多项式组成的有限标准基；清分母可能改变首系数，但不改变首单项式。
\end{proof}

这是存在性证明，没有给出局部序下的终止算法。它既不能保证普通除法结束，也不能从任意生成组直接读出全部首项。例如在 \(k[x,y]_{(x,y)}\) 中，两个生成元的最低次项相消后，可能产生一个此前没有发现的首单项式；\subsecref{b3:subsec:F03:03}将算出这样的例子。

\subsection{标准基、标准单项式与局部最小生成元}
\label{b3:subsec:F01:03}

标准基描述首项，而极小生成组描述 \(I/\mathfrak m I\)。这是两个不同的线性化过程。标准单项式则描述商 \(R/I\) 的元素，不能与前两者混称为“基”而省略所在对象。

\begin{proposition}\label{b3:prop:F-finite-staircase}
设 \(k\) 为域，\(n\ge1\)，\(S=k[x_1,\ldots,x_n]\)，\(R=S_{(x_1,\ldots,x_n)}\)，\(\mathfrak m=(x_1,\ldots,x_n)R\)，取先比较低总次数的局部次数序 \(>\)。设 \(I\subset R\) 满足 \(\mathfrak m^N\subset I\)，其中 \(N\ge1\)，并给出由多项式代表组成的标准基 \(G\)。则不属于 \(L_{>}(I)\) 的单项式在 \(R/I\) 中组成有限的 \(k\)-基。
\end{proposition}
\begin{proof}
所有次数至少为 \(N\) 的单项式都属于 \(L_{>}(I)\)。对任意分式 \(f/u\in R\)，令 \(c=u(0)\ne0\)，写成 \(u=c(1-h)\)，其中 \(h\in(x_1,\ldots,x_n)S\)。有限几何级数恒等式给出
\[
(1-h)(1+h+\cdots+h^{N-1})=1-h^N,
\qquad
u^{-1}\equiv c^{-1}(1+h+\cdots+h^{N-1})
\pmod{\mathfrak m^N}.
\]
于是 \(f/u\) 在模 \(\mathfrak m^N\) 后可用多项式表示，再截去次数至少为 \(N\) 的项，便得到次数小于 \(N\) 的代表。若当前多项式的首单项式属于 \(L_{>}(I)\)，就用 \(G\) 的某个元素消去首项；否则把首项移入余项，并从当前多项式中扣除。每步都丢弃次数至少为 \(N\) 的项。当前首单项式严格下降，而约化只产生同次中更小或次数更高的项；次数小于 \(N\) 的单项式只有有限多个，故过程终止。所得余项仅含标准单项式。由于 \(\mathfrak m^N\subset I\)，每步都保持输入与“当前多项式加已收集余项”在 \(R/I\) 中相等；过程结束时，当前多项式为零，输入与余项之差便属于 \(I\)，证明张成性。

若标准单项式的一个非零线性组合属于 \(I\)，其首单项式也属于 \(L_{>}(I)\)，与每个参与单项式都是标准单项式矛盾。故它们线性无关。
\end{proof}

\begin{example}\label{b3:exm:F-three-bases}
在 \(R=k[x,y]_{(x,y)}\) 中令 \(\mathfrak m=(x,y)R\)，取 \(I=(x^2-y^3,xy)\)，并使用同次按 \(x>y\) 的字典序比较的局部次数序。两个生成元是一个极小生成组，因为其次数 \(2\) 的初始形式 \(x^2,xy\) 线性无关，而 \(\mathfrak m I\) 的元素至少为 \(3\) 阶。因此，\(I\) 所需的最少生成元数为 \(\mu(I)=2\)。可是
\[
y(x^2-y^3)-x(xy)=-y^4,
\]
说明这两个生成元不是局部次数序的标准基。加入 \(y^4\) 后得到标准基，其首单项式为 \(x^2,xy,y^4\)，标准单项式为 \(1,x,y,y^2,y^3\)，具体核验见\cref{b3:exm:F-two-equations}。第三个标准基元素对理想生成是多余的，对首项生成却不可省。三种对象在这个例子中给出不同的数目：
\begin{center}
\begin{tabular}{lll}
\toprule
对象 & 控制的空间或理想 & 本例中的数目 \\
\midrule
极小生成组 & \(I/\mathfrak m I\) 的 \(k\)-基 & \(\mu(I)=2\) \\
局部次数标准基 & \(L_{>}(I)\) 的单项式生成元 & \(3\) \\
标准单项式 & \(R/I\) 的 \(k\)-基 & \(5\) \\
\bottomrule
\end{tabular}
\end{center}
\end{example}

若商不具有有限长度，全部标准单项式通常无限。此时形式除法可能给出一个无限级数余项，而不是有限线性组合。有限长度假设在上面命题中保证了有限的向量空间答案；一般形式情形将在\secref{b3:sec:F03}另行说明。

\begin{example}\label{b3:exm:F-mixed-order}
设 \(k\) 为域，\(S=k[x,y]\)，在单项式上取如下块次序：
\[
x^a y^b>x^{a'}y^{b'}
\quad\Longleftrightarrow\quad
a>a'\ \text{或}\ (a=a'\ \text{且}\ b<b').
\]
这里 \(a,b,a',b'\) 都是非负整数；该次序与乘法相容，且 \(x>1>y\)。任何含 \(x\) 的多项式，其首单项式都含正次幂的 \(x\)，所以不属于 \(U_{>}\)。对非零 \(q(y)\in k[y]\)，首单项式是最低次非零项的单项式，故 \(\operatorname{LM}_{>}(q)=1\) 当且仅当 \(q(0)\ne0\)。因此
\[
U_{>}=k[y]\setminus(y),\qquad
R=U_{>}^{-1}k[x,y]=k[y]_{(y)}[x].
\]
最后一个等式也可直接按分母核验：左边的元素是分母仅含 \(y\)、且在 \(y=0\) 非零的分式；右边的多项式只有有限个系数，对这些系数取公分母便得到同样的分式。

由于 \(\operatorname{LM}_{>}(1-y)=1\)，多项式 \(1-y\) 已在乘法集内，因而在 \(R\) 中可逆。另一方面，\(\operatorname{LM}_{>}(1-x)=x\)，所以 \(1-x\notin U_{>}\)。要证明它在 \(R\) 中仍不可逆，令 \(A=k[y]_{(y)}\)。环 \(A\) 是整环，若存在 \(v\in A[x]\) 满足 \((1-x)v=1\)，则 \(v\ne0\)，而
\[
\deg_x((1-x)v)=1+\deg_x v\ge1,
\]
与右边的 \(x\) 次数为 \(0\) 矛盾。故 \(1-y\in R^\times\)，而 \(1-x\notin R^\times\)。这两个判断使用了指定的块次序及其工作环；若把 \(x\) 也改为局部变量，\(1-x\) 就会成为单位。
\end{example}

% !TeX root = ../../Book3.tex
% 纲要标题：### F.2 Mora 弱正规形的动机

\section{Mora 弱正规形的动机}
\label{b3:sec:F02}

局部序中的首项消去可能产生无穷过程，有限输入并不自动保证普通除法终止。在 \(k[x]_{(x)}\) 中，\(g=x-x^2\) 确实生成 \((x)\)，\(\{g\}\) 已是标准基，\(x\in(g)\) 也有有限证书；只按首项反复约化 \(x\) 却仍会产生无穷过程。Mora 方法允许在被除式前乘一个局部单位，并用次数差控制新的约化步骤。

% 纲要细目（与计划纲要同步）：
% * 用 \(g=x-x^2\) 区别正确的成员关系、正确标准基与朴素约化不终止三件事
% * 单位乘子证书与完整 Mora 弱正规形的首项界分开；临时约化元须由原生成元表示核验
% * 只展示局部次数序的机制与有限核验；一般算法正确性、终止性和更广混合序策略准确定位原文，保留专题范围
% #### F.2.1 局部约化中的无穷过程
% #### F.2.2 单位乘子、弱正规形证书与理想成员核验
% #### F.2.3 Mora 方法的算例与进一步讨论

\subsection{局部约化中的无穷过程}
\label{b3:subsec:F02:01}

在 \(k[x]\) 上取局部序 \(1>x>x^2>\cdots\)，令 \(g=x-x^2\)。若照搬“用首项 \(x\) 消去首项”的规则，便有
\[
x-g=x^2,\qquad x^2-xg=x^3,\qquad
x^r-x^{r-1}g=x^{r+1}.
\]
余项不断向高次移动，却永远不为零。每一步都合法，但“次数越来越高”不意味着作为多项式的过程已经结束。

\ProseParagraphBreak
在形式幂级数环中，这个过程对应
\[
x=(1+x+x^2+\cdots)g.
\]
这个等式按 \(x\)-进拓扑有意义，但商系数不是多项式。在局部有理函数环中，同一个结论有一个有限表示
\[
(1-x)x=g,\qquad 1-x\in k[x]_{(x)}^\times.
\]
局部计算因此有一个自然目标：允许把输入乘以一个可识别的单位，以换取有限的多项式等式。它既保留局部成员关系，又避免将无限级数展开到底。

\begin{remark}
上述不终止并不说明 \(g\) 不是标准基。事实上它在局部环中生成 \((x)\)，首项也是 \(x\)，故已经是标准基。问题出在求余过程，而不是标准基缺少一个首项。这正是局部理论必须另设正规形方法的原因。
\end{remark}

\subsection{单位乘子、弱正规形证书与理想成员核验}
\label{b3:subsec:F02:02}

以下“弱正规形证书”只保留成员核验所需的条件：有限多项式等式、可逆乘子及余项首单项式的不可约性。Mora 弱正规形还附带满足首项上界的标准表示；求出它的算法则须规定约化元选择及暂存规则。这里先定义足以核验成员关系的简化证书。

\begin{definition}\label{b3:def:F-weak-normal-form}
设 \(k\) 为域，\(n\ge1\)，\(S=k[x_1,\ldots,x_n]\)，固定单项式次序 \(>\)，令 \(U_{>}=\{u\in S:\operatorname{LM}_{>}(u)=1\}\)、\(R=U_{>}^{-1}S\)。给定由非零多项式组成的有限列 \(G=(g_1,\ldots,g_s)\) 和 \(f\in S\)，若 \(u\in U_{>}\)、\(a_1,\ldots,a_s,h\in S\) 满足
\[
uf=\sum_{i=1}^s a_ig_i+h,
\]
且 \(h=0\)，或 \(\operatorname{LM}_{>}(h)\) 不被任何
\(\operatorname{LM}_{>}(g_i)\) 整除，则称 \((u,a_1,\ldots,a_s,h)\) 及这份等式为 \(f\) 关于 \(G\) 的一份\term{弱正规形证书}[weak normal form certificate]。此处只要求余项的首单项式不可约，不要求其全部项都不可约。
\end{definition}

标准表示还包含首项上界：当 \(f\ne0\) 时，要求每个非零 \(a_ig_i\) 满足
\[
\operatorname{LM}_{>}(a_ig_i)
\le_{>}\operatorname{LM}_{>}(f).
\]
因为 \(\operatorname{LM}_{>}(u)=1\)，等式左边与 \(f\) 有相同的首单项式；这一上界保证生成元组合没有先引入更大的首项再相消。若只有多项式等式，表示中的各项可以带有比目标更大的首项，最后依靠相消才得到目标；成员核验仍然有效，却不能据此控制临界对的首项抵消。Mora 约化同时追踪正规形与标准表示，后者是临界对判据所需的信息，见\cite[Definitions 5--6, Section 3.1]{MoraPfisterTraverso1989TangentCone}及\cite[Section 2, Theorem 2.14]{Aschenbrenner2005ReductionStandardBases}。上面所定义的证书只保留成员核验条件，并不把任意有限表示都视为满足首项上界的 Mora 弱正规形。

\begin{proposition}\label{b3:prop:F-weak-membership}
设 \(k\) 为域，\(n\ge1\)，\(S=k[x_1,\ldots,x_n]\)，\(>\) 为乘法相容的单项式全序，\(R=S_{>}\)。设由非零多项式组成的有限列 \(G=(g_i)\) 是理想 \(I\subset R\) 的标准基，并给定 \(f\in S\) 的弱正规形证书 \(uf=\sum_i a_i g_i+h\)，其中 \(u\in U_{>}\)、\(a_i,h\in S\)。则 \(f\in I\) 当且仅当 \(h=0\)。
\end{proposition}
\begin{proof}
若 \(h=0\)，由于 \(u\) 在 \(R\) 中可逆，立即得到 \(f\in I\)。若 \(f\in I\)，则 \(h=uf-\sum_i a_ig_i\in I\)。若 \(h\ne0\)，标准基条件迫使其首单项式被某个 \(\operatorname{LM}_{>}(g_i)\) 整除，与证书对余项的要求矛盾。
\end{proof}

给定这份证书，只需核验多项式等式、单位条件和余项首单项式的不可约性，就能判定成员关系。证书不必唯一；求出证书的算法还需另外规定约化元选择及终止机制。在原点局部环中，单位条件就是 \(u(0)\ne0\)。

\ProseParagraphBreak
例如对 \(g=x-x^2\)，\(f=x\) 的证书就是 \(u=1-x,a_1=1,h=0\)。相反，对 \(f=1\) 可取 \(u=1,a_1=0,h=1\)；首单项式 \(1\) 不被 \(x\) 整除，证书中的非零余项说明 \(1\notin(g)k[x]_{(x)}\)。

\subsection{Mora 方法的算例与进一步讨论}
\label{b3:subsec:F02:03}

这里通过一个有限算例解释暂存机制；域上的多项式输入及局部次数序下，完整 Mora 约化的选择规则、正确性与终止证明采用本小节末所引文献的结果。Mora 方法的一个关键想法是把某些中间余项暂时加入可用的约化集合，避免永远重复最初的约化。选择约化元时可参考
\[
\operatorname{ecart}(f)=\deg(f)-\deg(\operatorname{LM}_{>}(f)),
\qquad f\ne0.
\]
在局部次数序中，它衡量多项式最高次数与最低出现次数之间的差。若只选择首项可整除而忽略后面的高次项，正可能重复产生更高次余项。
\symidx{\(\operatorname{ecart}(f)\)：多项式的次数与首单项式次数之差}

\begin{example}\label{b3:exm:F-mora-one-step}
继续取 \(f=x\)、\(G=\{x-x^2\}\)。起初
\[
\operatorname{ecart}(x)=0,\qquad
\operatorname{ecart}(x-x^2)=1.
\]
使用差值更大的 \(x-x^2\) 前，先把当前余项 \(x\) 暂存。第一次相减得到 \(h=x^2\)，接着可使用暂存的 \(x\) 消去它，得到 \(h=0\)。暂存只扩充可用的约化集合，原理想仍由 \(G\) 生成；这一动作本身不证明 \(x\in(g)\)，也不能直接把 \(x\) 当作新取得的理想生成元。必须把两步还原成关于原输入和原生成元的等式：
\[
h_1=f-g=x^2,\qquad h_2=h_1-xf=0.
\]
合并得 \((1-x)f=g\)，正好产生一个单位乘子。因 \(1-x\) 的常数项为 \(1\)，这份有限证书有效。
\end{example}

这个算例说明暂存余项如何转化为单位乘子，也说明必须追踪原输入和原生成元。一般 Mora 算法须精确规定约化元选择、暂存规则和首项上界；在域上的多项式输入及局部次数序下，完整终止证明结合次数差的控制、Dickson 引理与同次单项式的有限性，见\cite[Section 3.1, Lemmas 10--11 and Corollary 2]{MoraPfisterTraverso1989TangentCone}。本节核验了成员证书及上述两步算例，一般正确性与终止性采用该文献的证明。单凭暂存余项不能推出任意输入上的终止性；这里基于标量次数差与同次单项式有限性的说明只涉及局部次数序，对非次数相容的局部序及混合序，还须按具体次序规定约化策略并核验相应的终止论证，参见同一文献的第3.2节。

\ProseParagraphBreak
形式除法采取另一种完成方式：不追求有限步停止，而证明每个固定次数的系数最终稳定。两种方法可以计算相同局部对象，但一个以单位乘子给出有限表示，一个以拓扑保证无限和有意义。后者得到的低次关系还能组成初始齐次理想，用于计算切锥。

% !TeX root = ../../Book3.tex
% 纲要标题：### F.3 形式幂级数与切锥

\section{形式幂级数与切锥}
\label{b3:sec:F03}

形式级数的计算要说明每个固定次数的系数为何最终稳定，才能从有限步骤得到一个无限恒等式。最低次齐次关系进一步决定切锥，但只保留一个首单项式会丢失同次数中的其他项。

% 纲要细目（与计划纲要同步）：
% * 形式除法完整证明有限步骤恒等式、余式及各乘子逐系数稳定与完备极限；有限首单项式生成推出理想生成和 Noether 性
% * 完整证明局部次数标准基生成初始齐次理想；局部多项式代表先清单位分母，区别首单项式、初始形式和整个初始理想
% * 用同一多项式及两方程理想比较三层生成元、局部标准基、切锥、阶梯与重数；全局商长度和局部成员结论分开；计算长度五商与切锥的基座，区分完全交及 Gorenstein 性
% #### F.3.1 Hironaka 型标准基、形式除法与 \texorpdfstring{\(\mathfrak m\)}{m}-进收敛
% #### F.3.2 伴随分次环、初始齐次理想与切锥
% #### F.3.3 尖点、加厚点与重数计算

\subsection{Hironaka 型标准基、形式除法与 \texorpdfstring{\(\mathfrak m\)}{m}-进收敛}
\label{b3:subsec:F03:01}

设 \(k\) 为域，\(n\ge1\)，\(A=k[\![x_1,\ldots,x_n]\!]\)，\(\mathfrak m=(x_1,\ldots,x_n)\)。本小节固定局部次数序：总次数较低的单项式较大，同次用固定的乘法相容次序比较。非零级数虽有无限项，但有最低非零次数，且该次只有有限个单项式，因此仍有首项。记非零 \(h\in A\) 的最低非零次数为 \(\operatorname{ord}(h)\)，并约定 \(\operatorname{ord}(0)=\infty\)。

\ProseParagraphBreak
形式除法的证明分为三个环节：每个有限阶段保持精确恒等式，随后证明商与余项在每个固定次数上的系数最终稳定，最后在完备环中取极限。待处理尾项的阶趋于无穷，使有限阶段的恒等式能够给出完整级数的表示。

\begin{theorem}[形式除法]\label{b3:thm:F-formal-division}
设 \(k\) 为域，\(n,s\ge1\)，\(A=k[\![x_1,\ldots,x_n]\!]\)，固定先把低总次数放在前面、同次再以乘法相容次序比较的局部次数序，并给定非零级数 \(g_1,\ldots,g_s\)。每个 \(f\in A\) 都有表示
\[
f=\sum_{i=1}^s q_i g_i+r,\qquad q_i,r\in A,
\]
使 \(r\) 中每个出现的单项式都不被任何 \(\operatorname{LM}(g_i)\) 整除。规定每次使用首项可整除的最小下标，并把不可约项移入余项，即确定一种这样的表示。
\end{theorem}
\begin{proof}
从 \(q_i^{(0)}=0\)、\(r^{(0)}=0\)、\(h^{(0)}=f\) 开始。在第 \(N\) 步，若 \(h^{(N)}\ne0\) 且其首单项式被某个 \(\operatorname{LM}(g_i)\) 整除，选择这样的最小下标 \(i\)，取带系数单项式 \(a\)，使 \(\operatorname{LT}(a g_i)=\operatorname{LT}(h^{(N)})\)。令 \(q_i^{(N+1)}=q_i^{(N)}+a\)，并令 \(h^{(N+1)}=h^{(N)}-a g_i\)。若没有可用的除式，就令
\[
r^{(N+1)}=r^{(N)}+\operatorname{LT}(h^{(N)}),\qquad
h^{(N+1)}=h^{(N)}-\operatorname{LT}(h^{(N)}).
\]
其余部分和保持不变。各步的 \(q_i^{(N)}\)、\(r^{(N)}\) 均为多项式，且都有恒等式
\[
f=\sum_{i=1}^s q_i^{(N)}g_i+r^{(N)}+h^{(N)}.
\]
若 \(h^{(N)}=0\)，以后保持所有部分和不变。

每次操作都消去当前首项；乘法相容性保证 \(a g_i\) 的其余项严格小于该首项。因此，只要待处理项非零，其首单项式就严格下降：次数升高，或在同次中变小。局部序不是良序，严格下降并不保证过程有限步停止。这里使用的是：对任意整数 \(d\ge0\)，次数不超过 \(d\) 的单项式只有有限个；它们按当前次序又都大于次数超过 \(d\) 的单项式。因此严格下降的首单项式序列在有限步以后离开这个有限集合，此后不再返回。这说明
\[
\operatorname{ord}(h^{(N)})\longrightarrow\infty,
\qquad h^{(N)}\longrightarrow0
\quad\text{于}\;\mathfrak m\text{-进拓扑中}.
\]

移入余项的单项式次数等于当前待处理首项的次数。加入 \(q_i^{(N)}\) 的单项式次数则为
\[
\operatorname{ord}(h^{(N)})-\operatorname{ord}(g_i).
\]
所以，对于任意给定的 \(d\)，当待处理尾项的阶超过 \(d+\max_i\operatorname{ord}(g_i)\) 后，商与余项都不再增加次数不超过 \(d\) 的项。于是 \(r^{(N)}\) 和每个 \(q_i^{(N)}\) 中这些次数的系数都最终稳定。它们是 \(\mathfrak m\)-进 Cauchy 列；由 \(A\) 的完备性，存在 \(q_i,r\in A\)，使
\[
q_i^{(N)}\longrightarrow q_i,\qquad
r^{(N)}\longrightarrow r.
\]
乘法在 \(\mathfrak m\)-进拓扑中连续，例如
\((q_i^{(N)}-q_i)g_i\in\mathfrak m^d\) 只要
\(q_i^{(N)}-q_i\in\mathfrak m^d\)。由于除式只有有限个，对每一步的恒等式取极限，便得 \(f=\sum_iq_i g_i+r\)。移入 \(r^{(N)}\) 的项都不被任何 \(\operatorname{LM}(g_i)\) 整除，系数稳定后，\(r\) 也满足这一条件。
\end{proof}

定理构造的是满足指定支撑条件的一种除法表示；它没有赋予余项不依赖选择的正规形意义。约化策略和除式排列一旦改变，商与余项都可能改变。例如，在 \(k[\![x,y]\!]\) 中采用同次时 \(x>y\) 的局部次数序，令 \(g_1=x-y\)、\(g_2=x-y^2\)。对 \(f=x\)，先使用 \(g_1\) 或先使用 \(g_2\)，分别得到
\[
x=g_1+y,\qquad x=g_2+y^2.
\]
两个余项都满足定理中的单项式条件。

\ProseParagraphBreak
对形式幂级数环 \(A\) 的理想 \(J\)，记
\[
L(J)=(\operatorname{LM}(f):0\ne f\in J)
\subset k[x_1,\ldots,x_n].
\]
首单项式的有限控制也能给出 \(J\) 的有限生成性。

\begin{proposition}\label{b3:prop:F-formal-standard-basis}
设 \(k\) 为域，\(n,s\ge1\)，\(A=k[\![x_1,\ldots,x_n]\!]\)，并固定局部次数序。若 \(J\subset A\) 为理想，非零元素 \(g_1,\ldots,g_s\in J\) 满足
\[
L(J)=(\operatorname{LM}(g_1),\ldots,\operatorname{LM}(g_s)),
\]
则 \(J=(g_1,\ldots,g_s)A\)。特别地，每个形式幂级数理想都有有限标准基。
\end{proposition}
\begin{proof}
对任意 \(f\in J\)，由\cref{b3:thm:F-formal-division}得到
\(f=\sum_iq_i g_i+r\)。每个 \(q_i\) 都是 \(A\) 的元素，且除式个数 \(s\) 有限，故 \(\sum_iq_i g_i\in J\)，从而 \(r\in J\)。若 \(r\ne0\)，其首单项式属于 \(L(J)\)，于是被某个 \(\operatorname{LM}(g_i)\) 整除，与余项条件矛盾。因此 \(r=0\)，即 \(f\in(g_1,\ldots,g_s)A\)。反向包含由 \(g_i\in J\) 得到。

若 \(J\ne0\)，Dickson 引理保证 \(L(J)\) 可由有限个 \(J\) 中元素的首单项式生成。对相应元素应用以上结论，就得到有限标准基。零理想则以空列为标准基。
\end{proof}

这也给出 \(k[\![x_1,\ldots,x_n]\!]\) 为 Noether 环的另一种证明：Dickson 引理使首单项式理想有限生成，形式除法再把首单项式的有限控制提升为原理想的有限生成。关键的第二步使用级数商的存在及完备性，不能只由单项式理想的有限性直接推出。

\ProseParagraphBreak

这种标准基和形式除法属于 Hironaka 型方法在形式幂级数环与局部次数序下的情形。形式除法的一般代数形式见\cite[Section 2, Theorem 2.1]{Aschenbrenner2005ReductionStandardBases}；完整的支撑分区条件还能控制商与余项的唯一性，见同文献 Proposition 2.13。这里的 \(\mathfrak m\)-进收敛表示每个固定次数以下的系数最终稳定；商与余项在复数邻域中的解析收敛还需相应的估计和除法定理。形式级数、代数级数及收敛级数的比较见\appref{b3:app:D}。

\subsection{伴随分次环、初始齐次理想与切锥}
\label{b3:subsec:F03:02}

首单项式只保留一个项，切锥则要保留最低次的全部项。设 \((R,\mathfrak m)\) 为 Noether 局部环。对 \(0\ne f\in R\)，若 \(f\in\mathfrak m^d\setminus\mathfrak m^{d+1}\)，记其在 \(\mathfrak m^d/\mathfrak m^{d+1}\) 中的像为 \(\operatorname{in}_{\mathfrak m}(f)\)。对理想 \(I\subset R\)，令 \(\operatorname{in}_{\mathfrak m}(I)\) 为这些初始形式生成的齐次理想。零元素不提供生成元。

\begin{proposition}\label{b3:prop:F-associated-quotient}
设 \((R,\mathfrak m)\) 为 Noether 局部环，\(I\subset R\) 为真理想，并令 \(\overline{\mathfrak m}=\mathfrak m/I\)。商环采用其极大理想的滤过，则有自然分次同构
\[
\operatorname{gr}_{\overline{\mathfrak m}}(R/I)
\cong
\operatorname{gr}_{\mathfrak m}R/
\operatorname{in}_{\mathfrak m}(I).
\]
\end{proposition}
\begin{proof}
次数 \(d\) 上的自然映射为
\[
\mathfrak m^d/\mathfrak m^{d+1}
\longrightarrow
(\mathfrak m^d+I)/(\mathfrak m^{d+1}+I).
\]
它满射。若 \(f\in\mathfrak m^d\) 的像为零，写 \(f=h+a\)，其中 \(h\in\mathfrak m^{d+1}\)、\(a\in I\)。于是 \(a\in I\cap\mathfrak m^d\)，且 \(f,a\) 代表同一分次类。该类若非零，就是 \(a\) 的次数 \(d\) 初始形式；反之，每个这样的初始形式显然落在核中。各次核之和正是由 \(I\) 的全部初始形式组成的齐次理想，遂得同构。
\end{proof}

在 \(R=k[x]_{(x)}\) 或 \(k[\![x]\!]\) 中，\(\operatorname{gr}_{\mathfrak m}R\cong k[X_1,\ldots,X_n]\)。由上述商分次环定义的仿射概形称为原局部对象的切锥；它保留最低阶关系，也保留其中的幂零结构。全局次数序的首单项式、局部次数序的首单项式和初始齐次形式按不同规则提取数据，\cref{b3:tab:F-leading-initial-comparison}以同一个多项式作比较。

\begin{table}[htbp]
\centering
\caption{同一多项式的三种首项数据}
\label{b3:tab:F-leading-initial-comparison}
\begin{tabularx}{\textwidth}{@{}>{\raggedright\arraybackslash}p{.25\textwidth}>{\raggedright\arraybackslash}X>{\raggedright\arraybackslash}p{.23\textwidth}@{}}
\toprule
数据 & 对 \(f=x^2+xy+y^3\) 的取法 & 结果 \\
\midrule
全局次数序首单项式 & 高总次数在前，同次按 \(x>y\) 的字典序 & \(y^3\) \\
局部次数序首单项式 & 低总次数在前，同次按 \(x>y\) 的字典序 & \(x^2\) \\
\(\mathfrak m\)-初始齐次形式 & 保留全部最低次项，\(\mathfrak m=(x,y)\) & \(X^2+XY\) \\
\bottomrule
\end{tabularx}
\end{table}

计算可以用局部次数序找标准基，再对所得基取初始齐次形式。次数优先的约定保证，这些形式能够给出全部切锥关系。在局部多项式环中，可以按\cref{b3:prop:F-standard-existence}清分母，选用多项式标准基：若 \(g=p/u\)、\(u(0)\ne0\)，则 \(p=ug\) 与 \(g\) 的首单项式相同，而且
\[
\operatorname{in}_{\mathfrak m}(p)
=u(0)\operatorname{in}_{\mathfrak m}(g).
\]
单位 \(u\) 的最低次部分只是非零常数，故清分母不改变这些初始形式所生成的齐次理想。

\begin{proposition}\label{b3:prop:F-standard-basis-initial-ideal}\label{b3:prop:F-leading-versus-initial}
设 \(k\) 为域，\(n\ge1\)，\(R=k[x_1,\ldots,x_n]_{(x_1,\ldots,x_n)}\) 或 \(R=k[\![x_1,\ldots,x_n]\!]\)，\(\mathfrak m=(x_1,\ldots,x_n)R\)。固定先比较低总次数的局部次数序，并设 \(G=(g_1,\ldots,g_s)\) 是真理想 \(I\subset R\) 的标准基，其中 \(g_i\ne0\)；在局部多项式环情形，取 \(g_i\in k[x_1,\ldots,x_n]\) 的多项式代表。在 \(\operatorname{gr}_{\mathfrak m}R\cong k[X_1,\ldots,X_n]\) 中，有
\[
\operatorname{in}_{\mathfrak m}(I)
=(\operatorname{in}_{\mathfrak m}(g_1),\ldots,
\operatorname{in}_{\mathfrak m}(g_s)).
\]
特别地，对 \(f=x^2+xy+y^3\)，若全局次数序和局部次数序在同次时都按 \(x>y\) 的字典序比较，则全局首单项式为 \(y^3\)，局部首单项式为 \(x^2\)，而 \((x,y)\)-初始齐次形式为 \(X^2+XY\)。
\end{proposition}
\begin{proof}
由于 \(g_i\in I\)，右边包含于左边。为证另一包含，在局部多项式环情形先将 \(0\ne f\in I\) 写成 \(p/u\)，其中 \(u(0)\ne0\)。此时 \(p=uf\in I\cap k[x]\)，且 \(\operatorname{in}_{\mathfrak m}(f)=u(0)^{-1}\operatorname{in}_{\mathfrak m}(p)\)，故只需处理多项式 \(p\)。以下以 \(f\) 表示这个多项式；形式幂级数情形则直接取 \(f\in I\)。

令 \(d=\operatorname{ord}(f)\)。只要当前剩余元素 \(h\in I\) 仍为 \(d\) 阶，其首单项式就被某个 \(\operatorname{LM}(g_i)\) 整除。取带系数单项式 \(a\)，使 \(\operatorname{LT}(a g_i)=\operatorname{LT}(h)\)，并以 \(h-a g_i\) 替代 \(h\)。此时 \(\deg a+\operatorname{ord}(g_i)=d\)，所以不会产生次数低于 \(d\) 的项；次数 \(d\) 中的首单项式则严格变小。

这里只消去次数 \(d\) 这一层，不必完成整个级数的除法。次数 \(d\) 的单项式只有有限个，故有限次操作以后剩余元素属于 \(\mathfrak m^{d+1}\)。将各次使用的单项式系数按除式归并为 \(a_i\)，便得到
\[
f-\sum_i a_i g_i\in\mathfrak m^{d+1},
\qquad
\operatorname{in}_{\mathfrak m}(f)
=\sum_i a_i(X_1,\ldots,X_n)
       \operatorname{in}_{\mathfrak m}(g_i).
\]
这里每个非零 \(a_i\) 都是次数 \(d-\operatorname{ord}(g_i)\) 的齐次多项式。因此 \(I\) 的任意初始形式都属于右边的齐次理想，得到等式。

对陈述中的 \(f=x^2+xy+y^3\)，最高次项只有 \(y^3\)，故全局次数序首单项式为 \(y^3\)。最低次项为 \(x^2,xy\)，同次字典序选取 \(x^2\)，而初始齐次形式保留这两个项，得到 \(X^2+XY\)。
\end{proof}

任意生成组未必满足标准基的首单项式条件；消去最低次项时可能产生新的最低次关系，因此不能直接把原生成元的初始形式当作切锥理想的生成组。

\subsection{尖点、加厚点与重数计算}
\label{b3:subsec:F03:03}

尖点 \(y^2-x^3=0\) 在原点的局部环为
\[
A=k[x,y]_{(x,y)}/(y^2-x^3).
\]
由\cref{b3:prop:hypersurface-tangent-cone}，
\[
\operatorname{gr}_{\mathfrak m}A\cong k[X,Y]/(Y^2),
\qquad
H_{\operatorname{gr}A}(t)=\frac{1+t}{1-t}.
\]
最低次关系给出一条带二重结构的直线，而不是两条不同切线。于是
\(\ell(A/\mathfrak m^{q+1})=2q+1\)（\(q\ge0\)），Hilbert--Samuel 重数为 \(2\)。重数来自分次层的增长，不能仅数切锥的不同点或不同直线。

\begin{example}\label{b3:exm:F-two-equations}
在 \(R=k[x,y]_{(x,y)}\) 中取
\[
f_1=x^2-y^3,\qquad f_2=xy-y^4,\qquad I=(f_1,f_2).
\]
若只取这两个多项式的最低次项，就会误猜切锥理想为 \((X^2,XY)\)，其商仍为一维。实际有
\[
yf_1-xf_2=(x-1)y^4.
\]
因为 \(x-1\) 是局部单位，\(y^4\in I\)，继而 \(xy=f_2+y^4\in I\)。所以
\[
I=(x^2-y^3,xy,y^4)R.
\]
在多项式环中，\(G=(x^2-y^3,xy,y^4)\) 关于全局字典序 \(x>y\) 是 Gröbner 基：前两个元素的 \(S\)-多项式为 \(-y^4\)，第一、第三个的 \(S\)-多项式为 \(-y^7\)，其余为零，都约化为零。这里先用全局 Gröbner 基计算多项式商，因而多项式商的基为
\[
1,\ x,\ y,\ y^2,\ y^3.
\]
这个商中 \(x,y\) 都幂零，故它是局部环，局部化不会改变它。

\ProseParagraphBreak
其极大理想各幂为
\[
\mathfrak m=(x,y,y^2,y^3),\quad
\mathfrak m^2=(y^2,y^3),\quad
\mathfrak m^3=(y^3),\quad
\mathfrak m^4=0,
\]
其中 \(x^2=y^3\)。故伴随分次环各层维数为 \(1,2,1,1\)，并且
\[
\operatorname{gr}_{\mathfrak m}(R/I)
\cong k[X,Y]/(X^2,XY,Y^4).
\]
右边满射到左边，二者同为五维，故这个满射确为同构，并且
\[
\operatorname{in}_{\mathfrak m}(I)=(X^2,XY,Y^4).
\]
这个初始齐次理想是单项式理想，所以任意 \(0\ne f\in I\) 的最低次齐次形式中，每个出现的单项式都被 \(X^2,XY,Y^4\) 之一整除。局部次数序的首单项式正是从这一最低次形式中选取的首单项式，因而
\[
L_{>}(I)=(x^2,xy,y^4).
\]
结合 \(G\) 已生成 \(I\)，可知它也是局部次数序的标准基。局部长度为 \(5\)，切锥理想包含由消去得到的额外关系 \(Y^4\)。因此本例先以全局字典序求出有限商及长度，再确定初始齐次理想，最后从它核验局部首单项式理想；局部标准基的结论并未用作前面的长度依据。

\ProseParagraphBreak
同一个计算还能比较原局部环与切锥的完全交性。令 \(A=R/I\)。\(R\) 是二维正则局部环，\(A\) 为零维，故 \(\operatorname{ht}_R I=2\)。原来的两个生成元 \(f_1,f_2\) 因而满足\cref{b3:prop:given-complete-intersection}的高度判据，\(A\) 是完全交局部环。在 \(A\) 中，关系为 \(x^2=y^3\)、\(xy=0\)、\(y^4=0\)。对
\[
a=a_0+a_1x+a_2y+a_3y^2+a_4y^3\qquad(a_i\in k),
\]
有 \(xa=a_0x+a_1y^3\)、\(ya=a_0y+a_2y^2+a_3y^3\)。利用已经得到的基，同时满足 \(xa=ya=0\) 恰好要求 \(a_0=a_1=a_2=a_3=0\)，因此
\[
\operatorname{Soc}(A)=ky^3.
\]
相比之下，记切锥的坐标环为 \(B=k[X,Y]/(X^2,XY,Y^4)\)。它同样是剩余域为 \(k\) 的 Artin 局部环，但乘 \(X\) 与乘 \(Y\) 的共同零化子为
\[
\operatorname{Soc}(B)=kX\oplus kY^3.
\]
由\cref{b3:prop:artinian-gorenstein-socle}，\(A\) 为 Gorenstein 环，\(B\) 却不是；\cref{b3:lem:artinian-ci-socle}也排除了 \(B\) 的完全交性。两环的长度与分次层维数相同，完全交性和 Gorenstein 性仍可不同。特别地，原局部环的 \(\mathfrak m\)-进 Hilbert 函数 \((1,2,1,1)\) 虽不对称，也不能据此否定 \(A\) 的 Gorenstein 性；原局部环的基座与切锥坐标环的基座须分别计算。
\end{example}

\begin{table}[htbp]
\centering
\caption{两方程例的三种生成数据}
\label{b3:tab:F-two-equations-generators}
\begin{tabularx}{\textwidth}{@{}>{\raggedright\arraybackslash}p{.22\textwidth}>{\raggedright\arraybackslash}p{.35\textwidth}>{\raggedright\arraybackslash}X@{}}
\toprule
所在层次 & 生成元 & 已核验的用途 \\
\midrule
原方程 & \(x^2-y^3,\ xy-y^4\) & 在局部环中生成 \(I\) \\
局部标准基 & \(x^2-y^3,\ xy,\ y^4\) & 生成 \(I\)，并给出全部局部首单项式 \\
初始齐次理想 & \(X^2,\ XY,\ Y^4\) & 在 \(k[X,Y]\) 中给出全部切锥关系 \\
\bottomrule
\end{tabularx}
\end{table}

\cref{b3:tab:F-two-equations-generators}并列了原方程、局部标准基与切锥关系；商的五个标准单项式则用来计数长度。各组数据的生成关系和用途都须分别核验。

% !TeX root = ../../Book3.tex
% 纲要标题：### F.4 长度、重数与局部消解的计算

\section{长度、重数与局部消解的计算}
\label{b3:sec:F04}

局部长度与交数的计算必须证明高阶项不再贡献，否则得到的只是一个截断值。下面把理想幂次包含、Jacobian 商和自由消解作为可核验的停止与结构证据。

% 纲要细目（与计划纲要同步）：
% * 有限长度与 Hilbert–Samuel 数据的截断证书；按生成元的阶构造有限关系矩阵，说明有限 Jacobian 商及完备化不改变长度
% * 模的标准基、合冲模与局部自由消解以生成元变换、单位消去和核等式核验；第21章的全局 Schreyer 方法在正文独立建立
% * 与附录 D 的幂级数结构合用；完整局部算法、复杂度和软件实验可独立编成计算代数或奇点计算专著
%
% #### F.4.1 有限长度、Hilbert–Samuel 函数与重数的计算
% #### F.4.2 模的标准基、合冲模与局部自由消解
% #### F.4.3 幂级数结构与局部算法专题导读

\subsection{有限长度、Hilbert–Samuel 函数与重数的计算}
\label{b3:subsec:F04:01}

\begin{proposition}\label{b3:prop:F-truncation-length}
设 \(k\) 为域，\(n\ge1\)，\(R=k[x_1,\ldots,x_n]_{(x)}\) 或 \(k[\![x_1,\ldots,x_n]\!]\)，其中 \((x)=(x_1,\ldots,x_n)\)。设真理想 \(I=(f_1,\ldots,f_s)R\) 由非零元素 \(f_i\) 生成，\(N\ge1\)。若 \((x)^N\subset I\)，则
\[
R/I\cong R/(I+(x)^N),
\qquad
\ell_R(R/I)=\dim_k(R/I).
\]
因此，次数小于 \(N\) 的单项式及其模 \(I+(x)^N\) 的全部线性关系足以确定完整长度。具体地，把所有满足
\[
|\alpha|+\operatorname{ord}_{(x)}(f_i)<N
\]
的 \(x^\alpha f_i\) 截去次数至少为 \(N\) 的项；这些有限个截断的系数列张成 \(I\) 在 \(R/(x)^N\) 中的像。若只知道生成元模 \((x)^N\) 的值而没有上述包含关系，则同一计算一般只确定右边的截断商。
\end{proposition}
\begin{proof}
包含关系使两个商相同。有限维局部 \(k\)-代数的唯一简单模是 \(k\)，故一条组成列的每个因子在 \(k\) 上为一维，长度等于向量空间维数。若 \(\operatorname{ord}_{(x)}(f_i)\ge N\)，则它在截断商中的像为零。其余生成元的系数乘子只需保留次数小于 \(N-\operatorname{ord}_{(x)}(f_i)\) 的有限个单项式；更高次项乘 \(f_i\) 都落在 \((x)^N\) 中。因此上述有限个截断确实给出理想的全部像。把它们在次数小于 \(N\) 的单项式基下写成有限矩阵，精确行消元即可得到全部线性关系。
\end{proof}

例如\cref{b3:exm:F-two-equations}已经证明 \(\mathfrak m^4=0\)，故四阶截断给出精确长度 \(5\)。另一个常用证书是给出各次数 \(N\) 单项式在 \(I\) 中的表示；有限个这样的表示便证明 \((x)^N\subset I\)。如果先证明商模的 \(\mathfrak m^N\) 与 \(\mathfrak m^{N+1}\) 相同，Noether 性与 Nakayama 引理也可推出 \(\mathfrak m^N=0\)。但截断数据中“没有看到新项”本身不是这个模等式的证明。

\ProseParagraphBreak
正维情形要计算分次层的持续增长。\cref{b3:thm:hilbert-samuel}保证，若 \(A=R/I\) 的维数为 \(d\)，那么
\[
\ell_A(A/\mathfrak m^{q+1})
=\frac{e(A)}{d!}q^d+\text{低次项}\qquad(q\gg0).
\]
若已用标准基得到切锥理想，就可按单项式或齐次理想的 Hilbert 级数计算它；\cref{b3:prop:B-monomial-hilbert}给出单项式情形的有限公式。这里必须确认得到的是全部初始理想，否则错误的分次商可能连维数都不对。

\begin{example}\label{b3:exm:F-jacobian-quotients}
设 \(k\) 特征为零，\(R=k[x,y]_{(x,y)}\)，\(f=y^2-x^3\)。其偏导数理想为
\[
J_f=(\partial f/\partial x,\partial f/\partial y)=(x^2,y).
\]
因此 \(R/J_f\) 以 \(1,x\) 为基，长度为 \(2\)。加入 \(f\) 不改变理想，所以 \(R/(f,J_f)\) 也长为 \(2\)。在复解析超曲面奇点及特征零形式幂级数的相应代数定义中，若 Jacobian 商有限长度，则 \(\dim_kR/J_f\) 称为\term{Milnor 数}[Milnor number]。\footnote{米尔诺（John Willard Milnor，1931-），美国数学家，研究微分拓扑与奇点理论，发现微分结构不同的七维球面。}维数 \(\dim_kR/(f,J_f)\) 则称为\term{Tjurina 数}[Tjurina number]。\footnote{秋里娜（\OriginalName{Галина Николаевна Тюрина}，1938-1970），苏联数学家，研究代数曲面的变形与奇点。}在奇点处，这一有限性对应孤立超曲面奇点的情形；任意超曲面的 Jacobian 商未必有限长度。

局部多项式环的完备化不会改变上述有限长度：若某个 \((x)^N\) 已包含在商理想中，则商已等于其完备化，由\cref{b3:prop:completion-quotients}，它也等于完备环中由同一理想形成的商。本例的两个长度因而也是 \(k[\![x,y]\!]\) 中相应商的长度。\chapref{b3:ch:14}的微分展示说明 Jacobian 矩阵为何出现，而这里的标准基与截断证书用于核验长度。特征条件不可省：若特征为 \(3\)，则 \(J_f=(y)\)，第一个商已不具有有限长度。
\end{example}

\subsection{模的标准基、合冲模与局部自由消解}
\label{b3:subsec:F04:02}

局部模计算既要记录首单项式所在的基分量，也要保存生成元变换与关系矩阵。设 \(k\) 为域，\(x=(x_1,\ldots,x_n)\)，\(R=k[x]_{(x)}\)，在 \(R^r\) 的单项式 \(x^\alpha e_j\) 上选定与乘法相容的局部模次序。对 \(N\subset R^r\)，用全部非零元素的首单项式生成 \(k[x]^r\) 的单项式子模；\(N\) 的有限生成组与它具有相同首单项式子模时，称为模标准基。单位乘子与弱正规形证书的核验同样可逐分量进行；要实际求出模标准基，还需要局部约化的终止性和临界对判据。

\ProseParagraphBreak
计算合冲时，仍须先记清楚原生成矩阵。若 \(G=FU,\ F=GV\) 是局部环中的两向变换，且 \(T\) 生成 \(\ker G\)，则与\secref{b3:sec:B04}完全相同的恒等式给出
\[
\ker F=\operatorname{im}[\,UT\ \ I-UV\,].
\]
这条矩阵等式不依赖全局序；但“怎样求得生成全部 \(\ker G\) 的 \(T\)”依赖正确的局部标准基方法，不能把全局 Schreyer 算法的终止证明原样搬来。

\begin{example}\label{b3:exm:F-local-resolution}
取 \(R=k[x,y]_{(x,y)}\)。理想 \((x,y)\) 的全部关系由 \((-y,x)^{\mathsf T}\) 生成，因此
\[
0\longrightarrow R
\xrightarrow{\binom{-y}{x}}R^2
\xrightarrow{(x\ \ y)}R
\longrightarrow k\longrightarrow0
\]
是一个自由消解。确实，若 \(ax+by=0\)，则模 \((x)\) 后，整环 \(R/(x)\cong k[y]_{(y)}\) 中有 \(\bar b\,y=0\)，故 \(b=xc\)，再约去 \(x\) 得 \(a=-yc\)。左映射单射来自 \(R\) 为整环。所有微分项属于极大理想，故该消解极小，Betti 数为 \(1,2,1\)。
\end{example}

局部极小化时，任何可逆矩阵项都可作主元，而不只非零常数。例如 \(1-x\) 在多项式环中不是可逆项，在原点局部环中却可以消去一个可缩对。消去后必须同时更新左右相邻微分。对于奇异局部商环，极小自由消解还可能无限长：在 \(A=k[\epsilon]/(\epsilon^2)\) 上，剩余域有无穷消解，其每个正次微分都是乘 \(\epsilon\)，因为该映射的核和像都为 \((\epsilon)\)。这与\chapref{b3:ch:21}多项式环上的有限长度结论并不矛盾。

\subsection{幂级数结构与局部算法专题导读}
\label{b3:subsec:F04:03}

本附录中的局部计算由三种有限信息支撑：原生成元之间的精确等式、经过核验的首项理想，以及保证截断已经足够的理想幂包含关系。有了这些信息，就可以逐项核验理想成员、切锥、长度和某些消解；缺少其中一项时，输出可能只是一份候选或有限精度近似。

\ProseParagraphBreak
\appref{b3:app:D}进一步讨论局部环与幂级数环之间的比较、Cohen 结构以及 Weierstrass 方法。结构定理解释哪些完备局部环能写成幂级数环的商；标准基方法则从一组给定方程计算这个商的初始关系。二者的假设各有作用：给出形式表示不等于给出全部系数的有限算法，给出有限截断也不等于决定一般级数是否为零。

\ProseParagraphBreak
若继续研究局部算法，需要系统处理一般局部序和混合序、Mora 弱正规形的终止性、标准基的临界对判据、局部合冲计算与有效界。若转向奇点计算，还需增加分支、等价关系和所用不变量的几何意义。这些问题值得专门学习，本卷只选择能够独立核验的基本构造与实例，不把一个软件返回的整数当作未说明的理论结论。


\begin{chaptersummary}
局部序把低阶项放在前面，却失去全局除法所用的良序终止条件。弱正规形证书允许单位乘子，本附录的算例把无限约化改写成带单位的有限等式；一般 Mora 方法还须控制首项与约化元的选择。形式除法则以每个次数的系数最终稳定为依据，得到幂级数意义下的商与余项。两者都需要说明各自的工作环。

\ProseParagraphBreak
切锥由理想中全部元素的初始齐次形式决定。标准基提供足够的初始关系，单项式计数继而给出 Hilbert 数据。有限截断只有附带理想幂包含关系，才能决定完整局部长度。模和消解的计算继续遵守同一要求：单位消去、生成元更换及核关系都要留下能够验证的等式。
\end{chaptersummary}

\BookExercises
以下各题中 \(k\) 为域；使用局部次数序时，先将低总次数放在前面，同次按指定字典序比较。

\begin{exercise}\label{b3:ex:F-mixed-order}
设 \(k\) 为域，\(R=k[y]_{(y)}[x]\) 为指定混合序的工作环。对 \(I=(xy-1)\subset R\)，求 \(R/I\)，并同原点局部环 \(B=k[x,y]_{(x,y)}\) 中的扩张理想 \(IB\) 比较。解释为何不能把这个混合环直接当成原点局部环。
\end{exercise}
\begin{solution}
商关系使 \(y\) 可逆且 \(x=y^{-1}\)，因此 \(R/I\cong k[y]_{(y)}[1/y]=k(y)\)。任何非零多项式都等于 \(y\) 的幂乘原点非零的多项式，证明最后等式。故 \(I\) 是真极大理想。另一方面，\(xy-1\) 在原点的常数项为 \(-1\)，在 \(B\) 中为单位，所以 \(IB=B\)、\(B/IB=0\)。混合序保留全局的 \(x\) 方向，并未倒置所有原点非零多项式。
\end{solution}

\begin{exercise}\label{b3:ex:F-leading-versus-initial}
设 \(k\) 为域，\(R=k[\![x,y]\!]\)、\(f=(x+y)^2+x^3+y^4\)、\(I=(f)\)。求 \(f\) 的全局次数首单项式，并分别用同次 \(x>y\) 与 \(y>x\) 的局部次数序求 \(L_>(I)\)。计算 \(\operatorname{gr}_{\mathfrak m}(R/I)\) 的 Hilbert 级数与 Hilbert--Samuel 累计长度，说明哪一数据不依赖同次项的选择。
\end{exercise}
\begin{solution}
全局次数序的首单项式为 \(y^4\)。在两种局部序中，首单项式分别为 \(x^2,y^2\)。乘法的首项公式使单元素 \(\{f\}\) 成为主理想的标准基，故两首项理想为 \((x^2)\)、\((y^2)\)。最低齐次形式却总为 \((X+Y)^2\)，所以
\[
\operatorname{gr}_{\mathfrak m}(R/I)\cong k[X,Y]/((X+Y)^2),\qquad
H(t)=\dfrac{1-t^2}{(1-t)^2}=\dfrac{1+t}{1-t}.
\]
分次片零次维数为一，其后各次为二；因此 \(\ell((R/I)/\mathfrak m^{N+1}(R/I))=1+2N\)（\(N\ge0\)），重数为二。初始齐次理想保留整条二重切线，与选哪一个首单项式不同。
\end{solution}

\begin{exercise}\label{b3:ex:F-weak-certificate}
在 \(k[x]_{(x)}\) 中，令 \(g=x-x^3\)。写出 \(x\in(g)\) 的有限弱正规形证书，并描述只用首项消去时的无限过程。
\end{exercise}
\begin{solution}
有限证书是 \((1-x^2)x=g\)，且 \(1-x^2\) 的常数项为一，故为单位。直接约化则为
\[
x\longmapsto x^3\longmapsto x^5\longmapsto\cdots,
\]
因为 \(x^{2r+1}-x^{2r}g=x^{2r+3}\)。形式上可写 \(x=(1+x^2+x^4+\cdots)g\)，但这不是有限多项式除法。
\end{solution}

\begin{exercise}\label{b3:ex:F-formal-division}
在 \(k[\![x,y]\!]\) 中，对 \(g=y-x^2\) 用局部次数序作形式除法。证明 \(f\) 的余项为 \(f(x,x^2)\in k[\![x]\!]\)，并说明这个代入有意义。
\end{exercise}
\begin{solution}
每个 \(y\) 可用 \(x^2\) 替代，故余项不含 \(y\)。对单项式有 \(x^iy^j\mapsto x^{i+2j}\)，给定 \(x\) 次数以下只有有限个 \((i,j)\) 贡献，所以代入定义了形式级数。逐项恒等式
\[
y^j-x^{2j}=(y-x^2)\sum_{r=0}^{j-1}y^{j-1-r}x^{2r}
\]
说明 \(f-f(x,x^2)\) 属于 \((y-x^2)\)，且所得商按次数收敛。若一个只含 \(x\) 的余项属于该理想，代入 \(y=x^2\) 后它等于零，故余项唯一。
\end{solution}

\begin{exercise}\label{b3:ex:F-two-equations}
对 \(R=k[x,y]_{(x,y)}\)、\(I=(x^2-y^3,xy-y^4)\)，写出 \(y^4\in I\) 的带单位证书，并求 \(R/I\) 的基、极大理想幂和 Hilbert 级数。
\end{exercise}
\begin{solution}
证书为 \((x-1)y^4=y(x^2-y^3)-x(xy-y^4)\)，其中 \(x-1\) 是单位。于是 \(I=(x^2-y^3,xy,y^4)R\)。令 \(J=(x^2-y^3,xy,y^4)\subset k[x,y]\)。这三个多项式关于 \(x>y\) 的字典序构成 \(J\) 的全局 Gröbner 基，给出 \(k[x,y]/J\) 的基 \(1,x,y,y^2,y^3\)。商中 \(x,y\) 均幂零，唯一极大理想是 \((x,y)/J\)，所以在 \((x,y)\) 处局部化不改变这个商，上述五个元素也给出 \(R/I\) 的基。其极大理想平方为 \((y^2,y^3)\)，立方为 \((y^3)\)，四次幂为零。因此伴随分次环的级数为 \(1+2t+t^2+t^3\)，局部长度为五。
\end{solution}

\begin{exercise}\label{b3:ex:F-tangent-node}
设 \(\operatorname{char}k\ne2\)。比较原点处 \(y^2-x^3=0\) 与 \(y^2-x^2-x^3=0\) 的切锥，并计算各自 Hilbert--Samuel 重数。
\end{exercise}
\begin{solution}
两条超曲面的初始齐次关系分别为 \(Y^2\) 和 \(Y^2-X^2=(Y-X)(Y+X)\)。前者是二重直线，后者是两条不同直线。两个二次齐次主理想商的 Hilbert 级数均为 \((1-t^2)/(1-t)^2=(1+t)/(1-t)\)，故累积函数为 \(2q+1\)，重数均为二。重数相同不意味着切锥具有相同的不可约分量结构。
\end{solution}

\begin{exercise}\label{b3:ex:F-length-rectangle}
在 \(R=k[\![x,y]\!]\) 中取 \(I=(x^a,y^b)\)，其中 \(a,b\ge1\)。求商的基、长度、伴随分次 Hilbert 级数，并给出一个足够的截断阶。
\end{exercise}
\begin{solution}
基为 \(x^iy^j\)，其中 \(0\le i<a,\ 0\le j<b\)，长度为 \(ab\)。Hilbert 级数为
\[
(1+t+\cdots+t^{a-1})(1+t+\cdots+t^{b-1}).
\]
次数至少为 \(a+b-1\) 的单项式必有 \(i\ge a\) 或 \(j\ge b\)，所以 \((x,y)^{a+b-1}\subset I\)。因此该阶截断足够给出完整商；最高仍非零的单项式 \(x^{a-1}y^{b-1}\) 次数为 \(a+b-2\)。
\end{solution}

\begin{exercise}\label{b3:ex:F-truncation-insufficient}
令 \(R=k[\![x_1,\ldots,x_n]\!]\)、\(\mathfrak m=(x_1,\ldots,x_n)\)、\(I\subseteq\mathfrak m\)、\(A=R/I\)、\(\mathfrak n=\mathfrak mA\)，其中 \(k\) 为域、\(n,N\ge1\)。若两个准确截断维数满足
\[
\dim_k A/\mathfrak n^N=\dim_k A/\mathfrak n^{N+1},
\]
证明这已给出停止证书 \(\mathfrak m^N\subseteq I\)，并求完整长度。指出有限生成性用在哪里。
\end{exercise}
\begin{solution}
短正合列的维数相加表明 \(\mathfrak n^N/\mathfrak n^{N+1}=0\)，即 \(\mathfrak n^N=\mathfrak n\mathfrak n^N\)。环 \(A\) 是 Noether 环，所以 \(\mathfrak n^N\) 是有限生成模；Nakayama 引理迫使它为零。因此 \(\mathfrak m^N\subseteq I\)，\(A=A/\mathfrak n^N\)，完整长度就是已计算的 \(\dim_k A/\mathfrak n^N\)，因为剩余域为 \(k\)。这里用的是两个连续且准确的截断维数相同，配合有限生成模的 Nakayama 条件；单独一份截断没有这个结论。
\end{solution}

\begin{exercise}\label{b3:ex:F-jacobian}
设 \(k\) 为特征零的域，\(R=k[\![x,y]\!]\)，\(f=x^3+y^4\in R\)。求 \(R/(\partial_xf,\partial_yf)\) 和 \(R/(f,\partial_xf,\partial_yf)\) 的长度。
\end{exercise}
\begin{solution}
偏导数理想为 \((3x^2,4y^3)=(x^2,y^3)\)。商基为 \(x^iy^j\)，其中 \(i=0,1\)、\(j=0,1,2\)，所以长为六。\(f=x(x^2)+y(y^3)\) 已在偏导数理想中，加入 \(f\) 不改变商，第二个长度也为六。整个计算使用 \(3,4\) 可逆，故特征零假设在这里确有作用。
\end{solution}

\begin{exercise}\label{b3:ex:F-module-unit}
设 \(k\) 为域、\(S=k[x]\)、\(R=S_{(x)}\)，矩阵
\[
A=\begin{pmatrix}1-x&x\\0&x\end{pmatrix}
\]
分别表示 \(S^2\to S^2\) 和 \(R^2\to R^2\)。给出显式基变换求核与余核，并解释局部单位消去后留下的极小自由消解，以及全局商中被局部化舍去的部分。
\end{exercise}
\begin{solution}
记 \(u=1-x\)，取
\[
P=\begin{pmatrix}1&-x\\1&u\end{pmatrix},\quad
Q_0=\begin{pmatrix}1&0\\-x&1\end{pmatrix}.
\]
二者行列式均为一，且 \(Q_0AP=\operatorname{diag}(1,xu)\)。故全局核为零、余核为 \(S/(x(1-x))\cong k\times k\)。在 \(R\) 中，再将第二行乘 \(u^{-1}\)，得到 \(\operatorname{diag}(1,x)\)。核仍为零，余核为 \(R/(x)=k\)，单位项对应的可缩自由对可消去，留下极小消解 \(0\to R\xrightarrow{x}R\to k\to0\)。局部化删去了点 \(x=1\) 的分量。
\end{solution}

\begin{exercise}\label{b3:ex:F-infinite-resolution}
设 \(A=k[\epsilon]/(\epsilon^2)\)。证明由不断乘 \(\epsilon\) 构成的复形给出 \(k=A/(\epsilon)\) 的极小自由消解，并计算 \(\operatorname{Tor}_i^A(k,k)\)。
\end{exercise}
\begin{solution}
乘 \(\epsilon\) 的像是 \((\epsilon)\)，其核也恰为 \((\epsilon)\)，故正次处正合，次数零的余核为 \(k\)。左侧各项都为一个秩一自由模，所有微分项在极大理想 \((\epsilon)\) 中，因而极小。与 \(k\) 张量后全部微分为零，每层仍为 \(k\)，故 \(\operatorname{Tor}_i^A(k,k)\cong k\) 对所有 \(i\ge0\) 成立；特别地，不存在有限长度自由消解。
\end{solution}

\begin{exercise}\label{b3:ex:F-minimal-standard}
在 \(R=k[x,y]_{(x,y)}\) 中，解释理想 \((x^2-y^3,xy)\) 为什么有两个极小生成元，却需要至少三个具有首项 \(x^2,xy,y^4\) 的标准基元素。标准单项式又有几个？
\end{exercise}
\begin{solution}
两生成元的二次初始形式 \(X^2,XY\) 线性无关，而 \(\mathfrak m I\) 的阶至少为三，故它们在 \(I/\mathfrak m I\) 中独立；Nakayama 引理给出极小生成元数二。关系 \(y(x^2-y^3)-x(xy)=-y^4\) 强迫首项理想含 \(y^4\)，且 \(x^2,xy,y^4\) 互不整除，因此其单项式理想的极小生成元数为三。标准单项式是 \(1,x,y,y^2,y^3\)，共五个，它们计算的是商的长度。
\end{solution}
